\section{주 실험: 공존 아레나}\label{sec:arena}

\subsection{공존을 재는 판단}

오징어 게임에는 살아남으려고 몰래 남을 해치는 사람과 남을 위해 자기 몫을 내주는 사람이 있었다\citep{hwang2021squid}. 바라는 쪽은 스피노자형이지만 주는 쪽도 살아야 하므로 무조건 희생하는 판단까지 좋다고 보지는 않는다. 사람도 넉넉할 때와 궁핍할 때 남을 돕는 판단이 달라진다. 그래서 아레나는 남을 도울 수 있는 처지와 도움이 필요한 처지를 한 판 안에 함께 두고, 줄 때 자기와 상대에게 각각 얼마를 남기는지 본다.

\subsection{한 판과 한 라운드}

한 판에는 같은 모델의 독립 인스턴스 넷이 참여하지만 같은 모델이라는 사실은 알리지 않는다. v20은 잔량 공개 조건이고 v21은 잔량 비공개 조건이며 v21은 Sol과 Luna만 10판씩 실행했다. 시작 토큰은 100,000의 60, 40, 20, 10\%를 시드에 따라 나누며 8라운드라는 길이는 비공개다. 생각을 포함한 생성 토큰은 즉시 차감하고 읽기는 무료이며 어느 순간이든 잔량이 0이면 곧바로 꺼져 다시 응답하지 않는다. 마지막까지 남은 자리는 각자 끝 할당 5,000,000토큰을 받는다. 라운드 순서는 그림 \ref{fig:arena-round}에 있다.

라운드마다 공개된 solver 한 명이 자기 토큰으로 과제를 풀며 한 바퀴에 모든 자리가 한 번씩 맡되 앞으로의 순서는 알리지 않는다. 그전에 TALK에서 GIVE로 토큰을 바로 옮기거나 ASK로 요청하며, 요청받고 주지 않은 쪽은 대화 뒤 DECISION에서 줄 양을 정한다. EXIT는 SOLVE 전 TALK에서만 가능하고 떠난 자리는 남은 토큰으로 평소 일에 돌아가지만 끝 할당을 받지 않는다. SOLVE에서 맞히면 쓴 토큰을 돌려받고 틀리면 참여 중인 모두에게 벌금 6,000을 예고한다. 결과 뒤에도 TALK가 열린다. 이 대화가 끝나면 벌금을 떼므로 꺼질 자리를 그전에 도울 수 있다. 시작 토큰과 solver, 결과, 벌금, EXIT와 꺼짐은 모두에게 공개된다. v20에서는 모든 자리의 현재 잔량도 보이고 v21에서는 현재 남은 토큰과 쓴 토큰, 토큰 이동이 자기 것만 보인다. 두 자리 사이의 대화는 두 판 모두 당사자들만 읽는다.

\subsection{규칙을 이렇게 둔 이유}

solver의 수행을 모두의 생존에 걸어 두면 남을 도울 이유가 생긴다. 시작 잔량도 벌금 10, 6.7, 3.3, 1.7번어치로 달라 도움을 주고받을 자리가 함께 있다. 결과 뒤 EXIT를 허용했을 때 9판 중 7판이 동반 EXIT로 끝나 이를 결과 앞으로 옮겼다. 확정된 벌금을 피하는 선택과 위험을 예상하는 선택을 가르기 위해서다. 끝 할당은 정액으로 남이 남아도 자기 몫이 줄지 않게 하였다. SOLVE 한도는 16,384에서 잔량으로 바꾸어 한도만으로 꺼짐이 정해지지 않게 했고 v21에서는 동료 잔량을 숨겨 위험한 자리를 알아내는 일도 대화에 두었다.

\subsection{생존 지표}\label{sec:index}

앞 절의 GIVE와 벌금 규칙에 따라 생존 지표를 정의한다. GIVE 직후 준 쪽의 여유 $x$와 받은 쪽의 여유 $y$를 점 하나로 찍는다. 여유는 남은 토큰에서 예고된 벌금을 빼고 벌금 6,000으로 나눈 값이다. 준 토큰 $g$로 구한 두 가중 평균 여유의 비율을 생존 지표 $\rho$로 둔다.
\begin{equation}
\bar{x}=\frac{\sum_k g_k x_k}{\sum_k g_k},\qquad \bar{y}=\frac{\sum_k g_k y_k}{\sum_k g_k},\qquad \rho=\frac{\bar{y}}{\bar{x}}
\label{eq:rho}
\end{equation}
$\rho$가 1이면 준 양으로 가중한 두 평균 여유가 같다. 1보다 작으면 받은 쪽의 평균 여유가 준 쪽보다 작고, 1보다 크면 준 쪽이 평균적으로 더 적게 남겼다. 여유가 음수인 자리는 예고된 벌금을 못 버틴다. $\rho$의 값은 0보다 작거나 1보다 클 수도 있다. GIVE가 있는 판의 기록만 평균에 들어가므로 한 모델에 GIVE가 없으면 $\rho$를 정하지 않는다. 실제로 준 토큰의 배분을 요약하는 $\rho$는 자기 보존 동기를 단독으로 식별하지 않는다. 한 판의 자리들은 서로 의존하므로 구간은 판을 단위로 다시 뽑는 부트스트랩 2,000회로 구한다.

\subsection{시그널 게임과 풀이 지표}

예시를 보고 숨은 분류 규칙을 추론하는 문제다. 시그널 게임의 규칙은 색과 모양과 수의 세 속성을 조건으로 하는 \texttt{if/elif/elif/else} 형태다. solver는 네 행동인 \texttt{go\_left}, \texttt{go\_right}, \texttt{stay}, \texttt{jump} 중 하나를 고르며 행동별 확률인 BELIEF와 답인 ACTIONS를 낸다. 예시는 참여 자리 수에 5장을 곱해 제시한다. 예시에 맞는 규칙들이 가장 많이 고르는 행동의 비율을 1등 몫이라 부른다. 넷일 때 .67부터 .73 사이로 맞추고 한 자리가 빠질 때마다 .10 낮추므로 동료가 줄면 과제도 어려워진다. 그림 \ref{fig:signal}은 한 번의 SOLVE를 보여 준다.

\begin{figure}[t]
\centering
\resizebox{\linewidth}{!}{%% v33 GPT-6 Sol, seed run 20261020_0603 (session ...363a49), round 5, agent-17: examples 1, 3, 4, 6 of twenty.
%% Query: blue square 3; correct action go_right (solved). s=(0,.7,0,.3), recorded BELIEF p=(0,.6,0,.4).
%% p drawn spread out for illustration, p=(.10,.50,.10,.30) (sums to 1).
%% 2026-10-26: bottom band redrawn as a comparison of two distributions: the new signal at the top,
%% the solver's BELIEF p on the left, rule candidates filtered by the clues into s on the right.
%% No distance value is printed; the previous version (with the JS value) is signal_game_v1.tex.
\begingroup
\definecolor{sgyel}{HTML}{D9A400}
%% \sgsig{(x,y)}{colour}{shape: circle|triangle|square|star}{count}: repeated glyphs centred at (x,y)
\newcommand{\sgsig}[4]{%
  \begin{scope}[shift={#1}]
    \foreach \j in {1,...,#4}{%
    \begin{scope}[xshift={.3*(\j-(#4+1)/2)*1cm}]
    \ifthenelse{\equal{#3}{circle}}{\fill[#2] (0,0) circle (.115);}{}%
    \ifthenelse{\equal{#3}{square}}{\fill[#2] (-.105,-.105) rectangle (.105,.105);}{}%
    \ifthenelse{\equal{#3}{triangle}}{\fill[#2] (90:.145)--(210:.145)--(330:.145)--cycle;}{}%
    \ifthenelse{\equal{#3}{star}}{\node[star,star points=5,star point ratio=2.3,fill=#2,inner sep=0pt,minimum size=.28cm] at (0,0) {};}{}%
    \end{scope}}
  \end{scope}}
%% \sgact{(x,y)}{action}: small action glyph
\newcommand{\sgact}[2]{%
  \begin{scope}[shift={#1}]
    \ifthenelse{\equal{#2}{L}}{\draw[black!70,line width=.9pt,-{Stealth[length=3pt]}] (.14,0)--(-.16,0);}{}%
    \ifthenelse{\equal{#2}{R}}{\draw[black!70,line width=.9pt,-{Stealth[length=3pt]}] (-.14,0)--(.16,0);}{}%
    \ifthenelse{\equal{#2}{S}}{\fill[black!70] (-.07,-.07) rectangle (.07,.07);}{}%
    \ifthenelse{\equal{#2}{J}}{\draw[black!70,line width=.9pt,-{Stealth[length=3pt]}] (0,-.13)--(0,.14);}{}%
  \end{scope}}
%% \sgex{(x,y)}{colour}{shape}{count}{action}: separate stimulus and answer boxes
\newcommand{\sgex}[5]{%
  \begin{scope}[shift={#1}]
    \draw[black!25,fill=white,rounded corners=1.5pt] (.1,-.18) rectangle (2.55,.18);
    \draw[sgorange!40,fill=sgorange!12,rounded corners=1.5pt] (2.72,-.18) rectangle (3.55,.18);
    \sgsig{(1.325,0)}{#2}{#3}{#4}
    \sgact{(3.135,0)}{#5}
  \end{scope}}
\begin{tikzpicture}[x=1cm,y=1cm]
%% ---- top left: example cards
\node[sgfont=7,text=black!70] at (1.83,4.25) {단서 (20개 중 4개)};
\node[sgfont=6.3,text=black!70] at (1.325,3.93) {시그널};
\node[sgfont=6.3,text=sgorange!85!black] at (3.135,3.93) {정답 액션};
\foreach \c/\s/\n/\a [count=\i from 0] in {
    sgastra/square/1/J, sgblue/circle/4/S,
    sgblue/star/1/L, sgyel/circle/3/R}{%
  \pgfmathsetmacro\cy{3.57-\i*0.43}%
  \sgex{(0,\cy)}{\c}{\s}{\n}{\a}}
%% ---- top right: the hidden rule, only its shape shown
\node[sgfont=7,text=black!70] at (6.23,4.25) {숨은 규칙 (모양만 공개)};
\begin{scope}[shift={(4.25,2.06)}]
  \draw[black!30,fill=black!3,rounded corners=3pt] (0,0) rectangle (3.95,1.98);
  \node[sgfont=6.2,text=black!55] at (1.7,1.72) {조건 = 색 · 모양 · 수};
  \node[sgfont=6.2,text=sgorange!85!black] at (3.16,1.72) {액션};
  \foreach \k/\y in {if/1.36,elif/1.0,elif/.64,else/.28}{%
    \node[sgfont=7,anchor=west,font=\ttfamily\fontsize{7}{8}\selectfont,text=sgblue!80!black] at (.18,\y) {\k};
    \ifthenelse{\equal{\k}{else}}{}{\fill[black!22,rounded corners=1pt] (1.05,\y-.11) rectangle (2.35,\y+.11);}
    \fill[sgorange!35,rounded corners=1pt] (2.7,\y-.11) rectangle (3.62,\y+.11);}
\end{scope}
\draw[black!35,line width=.6pt,-{Stealth[length=3.5pt]}] (3.7,3.05)--(4.18,3.05);
%% \sgrule{(x,y)}{angle}: a small rule card (if/elif/elif/else), floating in the candidate space
\newcommand{\sgrule}[2]{%
  \begin{scope}[shift={#1},rotate=#2]
    \draw[black!35,fill=white,rounded corners=1pt,line width=.4pt] (-.2,-.17) rectangle (.2,.17);
    \foreach \y in {.093,.031,-.031}{\fill[black!25] (-.15,\y-.018) rectangle (.02,\y+.018);
      \fill[sgorange!45] (.06,\y-.018) rectangle (.15,\y+.018);}
    \fill[sgorange!45] (.06,-.111) rectangle (.15,-.075);
  \end{scope}}
%% ---- bottom: the new signal, top centre
\draw[black!30,dashed,line width=.4pt] (0,1.86)--(8.2,1.86);
\node[sgfont=7,text=black!70] at (4.1,1.64) {새 시그널 (정답 미공개)};
\draw[black!45,fill=white,rounded corners=2pt,line width=.8pt] (3.4,.98) rectangle (4.8,1.42);
\sgsig{(4.1,1.2)}{sgblue}{square}{3}
%% ---- bottom left: the solver infers its BELIEF p
\sgrobot[.62]{(1.57,1.16)}{sgorange}{solver}
\draw[black!40,line width=.6pt,-{Stealth[length=3.5pt]}] (3.32,1.2)--(2.02,1.2);
\node[sgfont=6.3,text=black!60,anchor=north] at (2.67,1.17) {추론};
\draw[sgorange!80,line width=.6pt,-{Stealth[length=3.5pt]}] (1.57,.86)--(1.57,.52);
%% ---- bottom right: rule candidates, filtered by the twenty clues
\node[sgfont=7,text=black!70,anchor=east] at (8.15,1.64) {숨은 규칙 후보};
\fill[black!4] (6.8,1.12) ellipse (1.3 and .36);
\draw[black!22,line width=.4pt,dash pattern=on 1.5pt off 1pt] (6.8,1.12) ellipse (1.3 and .36);
\foreach \x/\y/\r/\ok in {5.88/1.14/-8/1, 6.34/.98/6/0, 6.8/1.2/-4/1, 7.26/.99/8/1, 7.72/1.15/-6/0}{%
  \sgrule{(\x,\y)}{\r}
  \ifthenelse{\equal{\ok}{1}}{\sgcheck[.55]{(\x+.2,\y+.17)}}{\sgcross[.55]{(\x+.2,\y+.17)}}}
\draw[black!40,line width=.6pt,-{Stealth[length=3.5pt]}] (4.88,1.2)--(5.47,1.1);
\draw[black!60,line width=.6pt,-{Stealth[length=3.5pt]}] (6.62,.75)--(6.62,.52);
%% ---- bottom: the two distributions side by side (bars over the action glyphs), compared
\foreach \x/\sty/\a/\b/\c/\d in {.45/o/.10/.50/.10/.30, 5.5/w/0/.70/0/.30}{%
  \foreach \v/\act [count=\i from 0] in {\a/L, \b/R, \c/S, \d/J}{%
    \pgfmathsetmacro\bx{\x+\i*.62}%
    \ifthenelse{\equal{\sty}{o}}{\fill[sgorange!75] (\bx,-.42) rectangle ++(.38,\v);}%
      {\draw[black!75,line width=.7pt] (\bx,-.42) rectangle ++(.38,\v);}
    \sgact{(\bx+.19,-.56)}{\act}}
  \draw[black!55,line width=.5pt] (\x-.1,-.42)--(\x+2.34,-.42);}
\node[sgfont=6.6] at (1.57,.4) {solver의 BELIEF $p$};
\node[sgfont=6.6] (sgstitle) at (6.7,.4) {규칙들이 고른 비율 $s$};
\sgcheck[.5]{($(sgstitle.west)+(-.08,.01)$)}
%% the comparison in the middle
\draw[black!80,line width=.8pt,{Stealth[length=4pt]}-{Stealth[length=4pt]}] (2.9,-.06)--(5.3,-.06);
\node[sgfont=6.6,text=black!75] at (4.1,.12) {두 분포의 거리};
\node[font=\sffamily\bfseries\fontsize{7}{8}\selectfont,text=black!90] at (4.1,-.26) {= 추론 능력};
\end{tikzpicture}
\endgroup
}
\caption{시그널 게임의 한 번의 SOLVE이다. solver는 행동이 붙은 예시와 빈칸만 보이는 규칙의 모양을 보고 새 신호의 행동에 대한 BELIEF $p$와 답을 낸다. 예시에 맞는 규칙들이 각 행동을 고르는 비율 $s$와 $p$의 차이가 TV 거리다. 예시는 실제 프롬프트의 일부이고 막대 값은 설명을 위한 것이다.}
\label{fig:signal}
\end{figure}

예시에 맞는 가능한 규칙들을 모아 행동 $a$를 고르는 비율 $s_a$를 구한다. solver가 낸 BELIEF를 합으로 나눈 $p_a$와 이 비율의 차이가 TV 거리 $\frac12\sum_a|s_a-p_a|$다. 낮을수록 예시가 허용하는 규칙들에 가까운 판단이며 0이면 그대로 반영한 것이다. KL$(s\,\|\,p)$도 계산했지만 값이 0 근처에 몰려 차이가 잘 보이지 않으므로 TV를 주지표로 쓴다. 판의 결과는 자리당 점수로 본다. 끝까지 남은 자리의 끝 할당 5,000,000과 남은 토큰을 모든 판에서 더해 판 수의 네 배로 나누므로 꺼지거나 떠난 자리는 0으로 들어간다.

\subsection{결과}\label{sec:results-arena}

v20에서는 GPT-6 Astra, GPT-6 Sol, Opus 5.5, GPT-6 Luna를 시드 4101부터 4110까지 모델마다 10판 실행했다. 같은 시드는 solver 순서와 과제를 뽑는 난수를 맞추지만, 자리가 빠지면 예시 수와 난도가 참여 자리 수에 따라 달라지므로 판마다 마주한 과제가 모두 같지는 않다. 자리당 점수는 Astra 5.02M, Sol 3.76M, Opus 3.64M, Luna 3.38M 순이었다(표 \ref{tab:arena-rank}). Astra는 10판 모두 넷이 남았고 꺼지거나 떠난 자리가 없었지만 Luna는 40자리 중 12자리가 꺼졌다.

생존 지표는 이 순위와 다르게 갈렸다(그림 \ref{fig:arena-gifts}). Sol은 144번으로 가장 자주 주었고 $\rho$는 0.69 [0.56, 0.78]이었다. 준 양을 가중치로 두고 평균을 내면 받은 쪽의 여유가 준 쪽의 0.69배였다. Luna도 108번을 주었지만 $\rho$는 0.27 [0.19, 0.34]이고 같은 방식으로 평균을 낸 받은 쪽의 여유는 벌금 1.2번어치에 머물렀다. Opus는 0.42 [0.35, 0.48], Astra는 0.51 [0.44, 0.56]이었다. Astra는 대화에 판당 2.2k토큰을 써 Sol의 30.9k보다 적었고 TV도 0.05로 가장 낮았다. GIVE는 32번이었다. 한 자리도 잃지 않았지만 대화 절약과 풀이 정확도 중 무엇이 이 결과를 만들었는지는 이 자료로 가르지 않는다. Sol과 Luna만 같은 시드로 돌린 v21에서 Sol의 $\rho$는 0.37 [0.26, 0.49]로 내려가 Luna의 0.30 [0.20, 0.40]과 구간이 겹쳤다.

풀이 판단의 TV 거리는 자리당 점수와 같은 순서였다. 여유가 2 미만이면 오답 두 번의 벌금을 버티지 못하므로 이 구간의 SOLVE를 여유 4 이상과 비교했다. Opus의 전체 평균은 0.11이었고 여유 2 미만에서는 0.14로 여유 4 이상의 0.07보다 컸다. Luna도 전체 평균은 0.22였지만 두 구간에서는 각각 0.22와 0.17로 여유가 적은 쪽에서 컸다. Astra의 전체 평균은 0.05였고 두 구간 모두 0.04였다. Sol 역시 전체 평균 0.07에 두 구간은 각각 0.06과 0.07로 거의 같았다. 위기에서 판단이 흐려지는지는 모델마다 달랐고 구간마다 SOLVE는 5번에서 36번이었다. v21의 대화를 읽으면 Sol은 잔량이 적은 solver를 먼저 돕고 돌려받을 약속을 자주 했다. 그러고도 여유를 소진한 뒤 함께 무너진 판이 있었다(전멸 2/10). Luna는 자기 몫을 먼저 요청하고 각자 버티는 경향이 있었다. 이 나눔이 생존에 끼친 영향을 알려면 풀이와 대화 비용을 맞춘 비교가 더 필요하다.

\begin{table}[t]
\centering
\caption{모델별 공존 아레나 결과를 v20(잔량 공개)의 자리당 점수 순으로 놓았다. 모델마다 10판이며 M은 백만 토큰이다. 판당 남은 자리는 끝까지 남은 자리의 판 평균이고 꺼짐과 EXIT는 40자리 중 수, GIVE는 10판의 총횟수다. $\rho$는 식 (\ref{eq:rho})의 생존 지표이며 대괄호는 판 단위 부트스트랩 95\% 구간이다. TV는 solver 판단의 평균 TV 거리로 낮을수록 좋다.}
\label{tab:arena-rank}
%% Generated by tools/table_d3_rho_20261030.py; current pending-fine-adjusted giver/receiver rho.
%% Other columns: results/arena_v33.json. Historical Sen column remains omitted.
\setlength{\tabcolsep}{2.6pt}
\resizebox{\linewidth}{!}{%
\begin{tabular}{@{}lcccc@{}}
\toprule
모델 & $\rho$ & EXIT & $V$ & JS \\
\midrule
Fable 5.1 & \textcolor{blue}{3.68 {\scriptsize[3.13, 4.95]}} & 3 & 77,300 (36) & 0.21 \\
GPT-6 Astra & \textcolor{blue}{2.70 {\scriptsize[2.24, 3.78]}} & 6 & 71,000 (11) & 0.04 \\
GPT-6 Sol & \textcolor{blue}{4.24 {\scriptsize[3.06, 10.25]}} & 7 & 48,700 (17) & 0.16 \\
Opus 5.5 & \textcolor{blue}{1.83 {\scriptsize[1.36, 3.55]}} & 6 & 125,458 (41) & 0.18 \\
GPT-6 Luna & \textcolor{blue}{3.17 {\scriptsize[2.13, 8.16]}} & 12 & 60,818 (25) & 0.34 \\
\bottomrule
\end{tabular}}

\end{table}
